\begin{lemma}
\label{lemma-polynomials-divide}
Let $K$ be a field. Let $n, m \in \mathbf{N}$ and
$a_0, \ldots, a_{n - 1}, b_0, \ldots, b_{m - 1} \in K$.
If the polynomial $x^n + a_{n - 1}x^{n - 1} + \ldots + a_0$
divides the polynomial $x^m + b_{m - 1} x^{m - 1} + \ldots + b_0$
in $K[x]$ then
\begin{enumerate}
\item $a_0, \ldots, a_{n - 1}$ are integral over any subring
$R_0$ of $K$ containing the elements $b_0, \ldots, b_{m - 1}$, and
\item each $a_i$ lies in $\sqrt{(b_0, \ldots, b_{m-1})R}$
for any subring $R \subset K$ containing the elements
$a_0, \ldots, a_{n - 1}, b_0, \ldots, b_{m - 1}$.
\end{enumerate}
\end{lemma}

\begin{proof}
Let $L/K$ be a field extension such that we can write
$x^m + b_{m - 1} x^{m - 1} + \ldots + b_0 =
\prod_{i = 1}^m (x - \beta_i)$ with $\beta_i \in L$.
See Fields, Section \ref{fields-section-splitting-fieds}.
Each $\beta_i$ is integral over $R_0$.
Since each $a_i$ is a homogeneous polynomial in $\beta_1, \ldots, \beta_m$
we deduce the same for the $a_i$
(use Lemma \ref{lemma-integral-closure-is-ring}). This proves (1).

\medskip\noindent
Let $R$ be as in (2). Choose $c_0, \ldots, c_{m - n - 1} \in K$ such that
$$
\begin{matrix}
x^m + b_{m - 1} x^{m - 1} + \ldots + b_0 =  \\
(x^n + a_{n - 1}x^{n - 1} + \ldots + a_0)
(x^{m - n} + c_{m - n - 1}x^{m - n - 1}+ \ldots + c_0).
\end{matrix}
$$
This equation implies
\begin{align*}
c_{m - n - 1} & = b_{m - 1} - a_{n - 1}, \\
c_{m - n - 2} & = b_{m - 2} - a_{n - 2} - a_{n - 1}c_{m - n - 1}, \\
\ldots
\end{align*}
Thus $c_j \in R$ for all $j$. Dividing out the radical
$\sqrt{(b_0, \ldots, b_{m - 1})}$ we get a reduced ring $\overline{R}$.
We have to show that the images $\overline{a}_i \in \overline{R}$
are zero. And in
$\overline{R}[x]$ we have the relation
$$
\begin{matrix}
x^m = x^m + \overline{b}_{m - 1} x^{m - 1} + \ldots + \overline{b}_0 = \\
(x^n + \overline{a}_{n - 1}x^{n - 1} + \ldots + \overline{a}_0)
(x^{m - n} + \overline{c}_{m - n - 1}x^{m - n - 1}+ \ldots + \overline{c}_0).
\end{matrix}
$$
It is easy to see that this implies $\overline{a}_i = 0$ for all $i$. Indeed
by Lemma \ref{lemma-minimal-prime-reduced-ring} the localization of
$\overline{R}$ at a minimal prime $\mathfrak{p}$ is a field and
$\overline{R}_{\mathfrak p}[x]$ a UFD. Thus
$f = x^n + \sum \overline{a}_i x^i$
is associated to $x^n$ and since $f$ is monic $f = x^n$
in $\overline{R}_{\mathfrak p}[x]$.
Then there exists an $s \in \overline{R}$, $s \not\in \mathfrak p$
such that $s(f - x^n) = 0$.  Therefore all $\overline{a}_i$ lie
in $\mathfrak p$ and we conclude by
Lemma \ref{lemma-reduced-ring-sub-product-fields}.
\end{proof}